%%%%%%%%%%%%%%%%%%%%%%%%%%%%%%%%%%%%%%%%%%%%%%%%%%%%%%%%%%%
% Vector Class
%%%%%%%%%%%%%%%%%%%%%%%%%%%%%%%%%%%%%%%%%%%%%%%%%%%%%%%%%%%
%%%%%%%%%%%%%%%%%%%
% Vector Creation %
%%%%%%%%%%%%%%%%%%%
% Create a new vector.
% Example: \textbackslash newVec\{myNewVec\}\{2, 3, 4\}
% Example: \textbackslash myNewVec \(\xrightarrow{} \begin{pmatrix}2\\3\\4\end{pmatrix}\)
% #1 out, var name to save the result (vector) to
% #2 vector, vector value that is assigned to #1
\newcommand\newVec[2]{%
    \expandafter\edef\csname #1\endcsname{\xintCSVtoList{#2}}%
}

% Create a new vector that is a multiple of pi.
% Example: \textbackslash newVec\{myNewVec\}\{0.5, 1, 2\}
% Example: \textbackslash myNewVec \(\xrightarrow{} \begin{pmatrix}\frac{1}{2}\pi\\\pi\\2\pi\end{pmatrix}\)
% #1 out, var name to save the result (vector) to
% #2 vector, vector value that is multiplied with pi and assigned to out
\newcommand\newVecPi[2]{%
    \def\tmpVec{\xintCSVtoList{#2}}
    \xintFor* ##1 in \tmpVec \do {%
        \newScalarPi{tmp}{##1}
        \xintifForFirst{%
            \newVec{#1}{\tmp}
        }{%
            \appendToVec{#1}{\csname #1\endcsname}{\tmp}
        }
    }
}

% Create a new vector variable with a random value between lower and
% upper bound. The value is obtained by creating a random value between 0 and 1
% and then interpolating between lower and upper bound.
% #1 out, var name to save the result (vector) to
% #2 lower bound, scalar that limits the possible values at the bottom
% #3 upper bound, scalar that limits the possible values at the top
% #4 TODO, TODO
\newcommand\randomVec[4]{
    \edef\innersequence{\xintSeq[+1]{+1}{#2}}%
    \xintFor* ##1 in \innersequence \do {%
        \FPrandom\tmp
        \lerpScalar{#3}{#4}{\tmp}
        \xintifForFirst{%
            \newVec{#1}{\lerpScalarResult}
        }{%
            \appendToVec{#1}{\csname #1\endcsname}{\lerpScalarResult}
        }
    }
}

% TODO
% Clone the given vector variable to a new variable
% \newcommand\cloneVec[2][cloneVecResult]{
% }

% Copy values from one vector variable to another.
% #1 out, var name to save the result (vector) to
% #2 vector, vector which values are assigned to out
% TODO does it make sence to have the name be an optional parameter?
\newcommand\copyVec[2][copyVecResult]{
    \xintFor* ##1 in #2 \do {%
        \xintifForFirst{%
            \newVec{#1}{##1}
        }{%
            \appendToVec{#1}{\csname #1\endcsname}{##1}
        }
    }
}

% TODO
% \newcommand\createVec[1]{
% }

% Set the values of a given vector.
% #1 out, var name to save the result (vector) to
% #2 vector, vector that is assigned to out
% \newcommand\setVec[2]{%
%     \newVector{#1}{#2}
% }

% TODO decide whether to use the low level assignments or use the functions like setVector to set the variable
% Set or create a vector with value 0.
% #1 out, var name to save the result (vector) to
% #2 dims, number of dimensions (rows)
\newcommand\zeroVec[2][3]{%
    \edef\sequence{\xintSeq[+1]{+1}{#1}}%
    \xintFor* ##1 in \sequence \do {%
        \xintifForFirst{%
            \newVec{#2}{0}
        }{%
            \appendToVec{#2}{\csname #2\endcsname}{0}
        }
    }
}

% TODO is it really necessary/useful to differentiate between #1 and #2
% Append a scalar to a already defined vector.
% #1 out, var name to save the result (vector) to
% #2 vector, the currently present values
% #3 scalar, the scalar to add to #2
\newcommand\appendToVec[3]{%
    \expandafter\edef\csname #1\endcsname{#2 \xintCSVtoList{#3}}%
}

%%%%%%%%%%%%%%%%%%%
% Vector Printing %
%%%%%%%%%%%%%%%%%%%
% Print the values of this vector variable in vector notation.
% Equation: \(\begin{pmatrix} v1\\ ...\\ vN \end{pmatrix}\)
% Example: \textbackslash newVec\{myNewVec\}\{1, 2, 3\}
% Example: \textbackslash printVec\{\textbackslash myNewVec\} \(\xrightarrow{} \begin{pmatrix}1\\2\\3\end{pmatrix}\)
% #1 fraction, whether or not the value should be printed as a fraction (val=1)
% #2 vector, the vector that should be printed
\newcommand\printVec[2][0]{%
    \ensuremath{%
        \begin{pmatrix}
            \xintFor* ##1 in {#2} \do {%
                \xintifForLast{
                    \FPifeq{#1}{0}
                        \roundScalar{##1}
                        \scalarRoundResult
                    \else
                        \numToFractionA{##1}
                    \fi
                }{%
                    \FPifeq{#1}{0}
                        \roundScalar{##1}
                        \scalarRoundResult
                    \else
                        \numToFractionA{##1}
                    \fi
                    \\
                }
            }
        \end{pmatrix}
    }
}

% Print the values of this vector variable in transposed notation.
% Equation: \(\trans{(v1, ..., vN)}\)
% Example: \textbackslash newVec\{myNewVec\}\{1, 2, 3\}
% Example: \textbackslash printVecT\{\textbackslash myNewVec\} \(\xrightarrow{} \trans{(1, 2, 3)}\)
% #1 fraction, whether or not the value should be printed as a fraction (val=1)
% #2 vector, the vector that should be printed
% TODO bug: last comma is shown when printing some fractions 0.5 and 0.25 work, 1.5, 2.5, 3.5 do not (all have number >1)
\newcommand\printVecT[2][0]{%
    \ensuremath{\trans{\left(
        \xintFor* ##1 in {#2} \do {%
            \xintifForLast{
                \FPifeq{#1}{0}
                    \roundScalar{##1}
                    \scalarRoundResult
                \else
                    \numToFractionA{##1}
                \fi
            }{%
                \FPifeq{#1}{0}
                    \roundScalar{##1}
                    \scalarRoundResult
                \else
                    \numToFractionA{##1}
                \fi
                ,
            }
        }
    \right)}}%
}

% Print the value of the given vector in point notation.
% Equation: \(\begin{bmatrix}#2\end{bmatrix}\)
% Example: \textbackslash newVec\{myNewVec\}\{1, 2, 3\}
% Example: \textbackslash printAsPoint\{\textbackslash myNewVec\} \(\xrightarrow{} \begin{bmatrix}1\\2\\3\end{bmatrix}\)
% #1 fraction, whether or not the value should be printed as a fraction (val=1)
% #2 vector, the vector that should be printed
\newcommand\printAsPoint[2][0]{%
    \ensuremath{%
        \begin{bmatrix}
            \xintFor* ##1 in {#2} \do {%
                \FPifeq{#1}{0}
                    \roundScalar{##1}
                    \scalarRoundResult
                \else
                    \numToFractionA{##1}
                \fi
                \xintifForLast{}{%
                    \\
                }
            }
        \end{bmatrix}
    }
}

% Print the value of the given vector as a multiple of Pi in vector notation.
% Equation: \(\begin{pmatrix}\left\{\frac{#2}{\pi}\right\}\end{pmatrix}\pi\)
% Example: \textbackslash newVec\{myNewVec\}\{0.5, 1, 2\}
% Example: \textbackslash printVecAsPi\{\textbackslash myNewVec\} \(\xrightarrow{} \begin{bmatrix}0.5\pi\\\pi\\2\pi\end{bmatrix}\)
% #1 fraction, whether or not the value should be printed as a fraction (val=1)
% #2 vector, the vector that should be printed
\newcommand\printVecAsPi[2][0]{%
    \ensuremath{%
        \begin{pmatrix}
            \xintFor* ##1 in #2 \do {%
                \printScalarAsPi[#1]{##1}
                % \FPifeq{#1}{0}
                %     \roundScalar{##1}
                %     \scalarRoundResult
                % \else
                %     \numToFractionA{##1}
                % \fi
                \xintifForLast{}{%
                    \\
                }
            }
        \end{pmatrix}
    }
}

% Print the value of the given vector as a multiple of Pi in transposed notation.
% Equation: \(\trans{\left(\left\{\frac{#2}{\pi}\right\}\right)}\pi\)
% Example: \textbackslash newVec\{myNewVec\}\{0.5, 1, 2\}
% Example: \textbackslash printVecAsPi\{\textbackslash myNewVec\} \(\xrightarrow{} \trans{(0.5\pi\\\pi\\2\pi)}\)
% #1 fraction, whether or not the value should be printed as a fraction (val=1)
% #2 vector, the vector that should be printed
\newcommand\printVecAsPiT[2][0]{%
    \ensuremath{\trans{\left(
        \xintFor* ##1 in #2 \do {%
            \printScalarAsPi[#1]{##1}
            \xintifForLast{}{%
                ,
            }
        }
    \right)}}
}

% TODO xintListWithSep{, }{#1}%
% Print all values of the given vector commaseparated.
% Equation: \(v1, ..., vN\)
% #1 vector, the vector which values should be printed
\newcommand\printVecContent[1]{%
    \xintNthElt{1}{#1},\xintNthElt{3}{#1},\xintNthElt{2}{#1}%
}

% Returns the x and y component of the given vector.
% Equation: \(X(v), Y(v)\)
% #1 vector, the vector which values should be printed
\newcommand\printVecContentXY[1]{%
    \xintNthElt{1}{#1}, \xintNthElt{2}{#1}
}

% Returns the x, y and z component of the given vector.
% Equation: \(X(v), Y(v), Z(v)\)
% #1 vector, the vector which values should be printed
\newcommand\printVecContentXYZ[1]{%
    \xintNthElt{1}{#1}, \xintNthElt{2}{#1}, \xintNthElt{3}{#1}%
}

% Returns the x, y, z and w component of the given vector
% Equation: \(X(v), Y(v), Z(v), W(v)\)
% #1 vector, the vector which values should be printed
\newcommand\printVecContentXYZW[1]{%
    \xintNthElt{1}{#1}, \xintNthElt{2}{#1}, \xintNthElt{3}{#1}, \xintNthElt{4}{#1}
}

% Returns the number of components that given vector has (num dims).
% #1 out, var name to save the result (scalar) to
% #2 vector, the vector of which to return the height
\newcommand\getVecHeight[2][vecHeightResult]{%
    \newScalar{#1}{\expandafter\xintLength\expandafter{#2}}
}

% Print the intermediate components of the scalarproduct calculation.
% Equation: \(D1(v1) * D1(v2) + ... + DN(v1) * DN(v2)\)
% Error: When both vectors do not have the same length
% #1 fraction, whether or not the value should be printed as a fraction (val=1)
% #2 vector, the first vector
% #3 vector, the second vector
% TODO this has a dangling + sign at the end when last elem is a fraction maybe fixable when copying code into the xintforlast
\newcommand\printScalarPComp[3][0] {%
    \FPifeq{\expandafter\xintLength\expandafter{#2}}{\expandafter\xintLength\expandafter{#3}}
    \else
        \PackageError{glmatrix package}{In scalar product}{vector A and B do not have same length}
    \fi
    \edef\innersequence{\xintSeq[+1]{+1}{\expandafter\xintLength\expandafter{#2}}}%
    \ensuremath{
        \xintFor* ##1 in \innersequence \do {%
            \FPifeq{#1}{0}
                \printScalarRounded{\xintNthElt{##1}{#2}}*\printScalarRounded{\xintNthElt{##1}{#3}}
            \else
                \numToFractionA{\xintNthElt{##1}{#2}}*\numToFractionA{\xintNthElt{##1}{#3}}
            \fi
            \xintifForLast{}{%
                +
            }
        }
    }
}

% Print the intermediate components of the crossproduct calculation.
% Equation: \(\begin{pmatrix} ab*bc - ac*bb\\ ac*ba-aa*bc\\ aa*bb-ab*ba \end{pmatrix}\)
% Error: When one of both vectors does not have length 3
% #1 fraction, whether or not the value should be printed as a fraction (val=1)
% #2 vector, the first vector
% #3 vector, the second vector
\newcommand\printCrossPComp[3][0] {%
    \FPifeq{\expandafter\xintLength\expandafter{#2}}{3}
    \else
        \PackageError{glmatrix package}{In cross product}{vector A is not of length 3}
    \fi
    \FPifeq{\expandafter\xintLength\expandafter{#3}}{3}
    \else
        \PackageError{glmatrix package}{In cross product}{vector A is not of length 3}
    \fi
    \FPset\aa{\xintNthElt{1}{#2}}
    \FPset\ab{\xintNthElt{2}{#2}}
    \FPset\ac{\xintNthElt{3}{#2}}
    \FPset\ba{\xintNthElt{1}{#3}}
    \FPset\bb{\xintNthElt{2}{#3}}
    \FPset\bc{\xintNthElt{3}{#3}}
    %
    \ensuremath{
        \FPifeq{#1}{0}
            \begin{pmatrix}
                \printScalarRounded{\ab}*\printScalarRounded{\bc} - \printScalarRounded{\ac}*\printScalarRounded{\bb}\\
                \printScalarRounded{\ac}*\printScalarRounded{\ba} - \printScalarRounded{\aa}*\printScalarRounded{\bc}\\
                \printScalarRounded{\aa}*\printScalarRounded{\bb} - \printScalarRounded{\ab}*\printScalarRounded{\ba}
            \end{pmatrix}
        \else
            \begin{pmatrix}
                \numToFractionA{\ab}*\numToFractionA{\bc} - \numToFractionA{\ac}*\numToFractionA{\bb}\\
                \numToFractionA{\ac}*\numToFractionA{\ba} - \numToFractionA{\aa}*\numToFractionA{\bc}\\
                \numToFractionA{\aa}*\numToFractionA{\bb} - \numToFractionA{\ab}*\numToFractionA{\ba}
            \end{pmatrix}
        \fi
    }
}

%%%%%%%%%%%%%%%%%%%%%%%
% Vector Modification %
%%%%%%%%%%%%%%%%%%%%%%%
% TODO add default param scalarRoundResult does not work since
% there is already a default param
% Round the values ot the given vector.
% #1 digits, number of digits after comma
% #2 out, var name to save the result (scalar) to
% #3 vector, the vector that schould be rounded
\newcommand\roundVec[3][4]{%
    \xintFor* ##1 in #3 \do {%
        \FPround\tmp{##1}{#1}
        \xintifForFirst{%
            \newVec{#2}{\tmp}
        }{%
            \appendToVec{#2}{\csname #2\endcsname}{\tmp}
        }
    }
}

% Ceil the values of thr given vector.
% #1 out, var name to save the result (vector) to
% #2 vector, the vector that should be ceiled
\newcommand\ceilVec[2][ceilVecResult]{
    \xintFor* ##1 in #2 \do {%
        \FPifint{##1}
        \else
            \FPadd\tmp{##1}{1}
        \fi
        \FPtrunc\tmp\tmp{0}
        \xintifForFirst{%
            \newVec{#1}{\tmp}
        }{%
            \appendToVec{#1}{\csname #1\endcsname}{\tmp}
        }
    }
}

% Floor the values of the given vector.
% #1 out, var name to save the result (vector) to
% #2 vector, the vector that should be floored
\newcommand\floorVec[2][floorVecResult]{
    \xintFor* ##1 in #2 \do {%
        \FPtrunc\tmp{##1}{0}
        \xintifForFirst{%
            \newVec{#1}{\tmp}
        }{%
            \appendToVec{#1}{\csname #1\endcsname}{\tmp}
        }
    }
}

% TODO should this be dehomogenVecResult?
% Dehomogenize the values of the given vector.
% #1 out, var name to save the result (vector) to
% #2 vector, the vector that should be dehomogenized
\newcommand\dehomogenVec[2][dehomogenResult]{
    % TODO check if 4 long or get last comp
    \FPset\tmpDiv{\xintNthElt{4}{#2}}
    % \FPadd\tmpDiv\tmpDiv{1}
    \edef\numElems{#2}
    \edef\vecLength{\xintSeq[+1]{+1}{\expandafter\xintLength\expandafter{\numElems}}}%
    \xintFor* ##1 in \vecLength \do {%
        \xintifForFirst{%
            \FPdiv\divResult{\xintNthElt{##1}{#2}}{\tmpDiv}
            \newVec{tmpVec}{\divResult}
        }{%
            % \xintifForLast{
            % }{
                \FPdiv\divResult{\xintNthElt{##1}{#2}}{\tmpDiv}
                \appendToVec{tmpVec}{\tmpVec}{\divResult}
            % }
        }
    }
    \copyVec[#1]{\tmpVec}
}

%%%%%%%%%%%%%%%
% Vector Math %
%%%%%%%%%%%%%%%
% Add the values of n vectors.
% Equation: \(out = vectors[1] + vectors[2] + ...\)
% Example: \textbackslash newVec\{myNewVecA\}\{0.5, 1, 2\}
% Example: \textbackslash newVec\{myNewVecB\}\{0.5, 1, 2\}
% Example: \textbackslash addVec\{\textbackslash myNewVecA\}\{\textbackslash myNewVecB\} \(\xrightarrow{} \begin{pmatrix}1\\2\\4\end{pmatrix}\)
% #1 out, var name to save the result (vector) to
% #2 vectors, list of vectors to sum
% TODO add check that the variable name that is given does not match any other parameter, since this does not work (inplace)
% TODO rewrite to have one vec #2 and a list of vecs #3 to add
\newcommand\addVec[2][addVecResult]{
    \def\veclist{\xintCSVtoList{#2}}
    \edef\elemOne{\xintNthElt{1}{\veclist}}
    \edef\innersequence{\xintSeq[+1]{+1}{\expandafter\xintLength\expandafter{\elemOne}}}%
    \xintFor* ##1 in \innersequence \do {%
        \FPset\sumCoord{0}%
        \xintifForFirst{%
            \xintFor* ##2 in \veclist \do {%
                \FPadd\sumCoord\sumCoord{\xintNthElt{##1}{##2}}%
            }
            \newVec{#1}{\sumCoord}
        }{%
            \xintFor* ##2 in \veclist \do {%
                \FPadd\sumCoord\sumCoord{\xintNthElt{##1}{##2}}%
            }
            \appendToVec{#1}{\csname #1\endcsname}{\sumCoord}
        }
    }
}

% Subract list of vectors from first vector of that list.
% Equation: \(out = vectors[1] - vectors[2] - ...\)
% Example: \textbackslash newVec\{myNewVecA\}\{0.5, 1, 2\}
% Example: \textbackslash newVec\{myNewVecB\}\{0.5, 1, 2\}
% Example: \textbackslash subVec\{\textbackslash myNewVecA\}\{\textbackslash myNewVecB\} \(\xrightarrow{} \begin{pmatrix}0\\0\\0\end{pmatrix}\)
% #1 out, var name to save the result (vector) to
% #2 vectors, list of vectors
% TODO add check that the variable name that is given does not match any other parameter, since this does not work
% TODO rewrite to have one vec #2 and a list of vecs #3 to add
\newcommand\subVec[2][subVecResult]{
    \def\veclist{\xintCSVtoList{#2}}
    \edef\elemOne{\xintNthElt{1}{\veclist}}
    \edef\innersequence{\xintSeq[+1]{+1}{\expandafter\xintLength\expandafter{\elemOne}}}%
    \xintFor* ##1 in \innersequence \do {%
        \FPset\subCoord{0}%
        \xintifForFirst{%
            \xintFor* ##2 in \veclist \do {%
                \xintifForFirst{%
                    \FPset\subCoord{\xintNthElt{##1}{##2}}%
                }{%
                    \FPsub\subCoord\subCoord{\xintNthElt{##1}{##2}}%
                }
            }
            \newVec{#1}{\subCoord}
        }{%
            \xintFor* ##2 in \veclist \do {%
                \xintifForFirst{%
                    \FPset\subCoord{\xintNthElt{##1}{##2}}%
                }{%
                    \FPsub\subCoord\subCoord{\xintNthElt{##1}{##2}}%
                }
            }
            \appendToVec{#1}{\csname #1\endcsname}{\subCoord}
        }
    }
}

% Multiply the values of a vector with these of a list of vectors. The list can have 1-N entries.
% Equation: \(out = vector * vectors[1] * vectors[2] * ...\)
% Example: \textbackslash newVec\{myNewVecA\}\{0.5, 1, 2\}
% Example: \textbackslash newVec\{myNewVecB\}\{0.5, 1, 2\}
% Example: \textbackslash mulVec\{\textbackslash myNewVecA\}\{\textbackslash myNewVecB\} \(\xrightarrow{} \begin{pmatrix}0.25\\1\\4\end{pmatrix}\)
% #1 out, var name to save the result (vector) to
% #2 vector, the vector to multiply
% #3 vectors, the list of vectors with which to multiply
\newcommand\mulVec[3][mulVecResult]{
    \def\vecList{\xintCSVtoList{#3}}
    \getVecHeight{#2}
    \edef\sequence{\xintSeq[+1]{+1}{\vecHeightResult}}%
    % Iterate each row / dimension
    \xintFor* ##1 in \sequence \do {%
        % Assign value of start vector (param #2)
        \FPset\mulValue{\xintNthElt{##1}{#2}}%
        % If this is the first iteration create a new vector
        \xintifForFirst{%
            % Iterate all vectors that are part of param #3
            \xintFor* ##2 in \vecList \do {%
                % Multiply with row value of current vector
                \FPmul\mulValue\mulValue{\xintNthElt{##1}{##2}}%
            }
            \newVec{#1}{\mulValue}
        }{%
            \xintFor* ##2 in \vecList \do {%
                \FPmul\mulValue\mulValue{\xintNthElt{##1}{##2}}%
            }
            \appendToVec{#1}{\csname #1\endcsname}{\mulValue}
        }
    }
}

% Divide the values of a vector with these of a list of vectors. The list can have 1-N entries.
% Equation: \(out = (vector / vectors[1] / vectors[2] / ...)\)
% Example: \textbackslash newVec\{myNewVecA\}\{0.5, 1, 2\}
% Example: \textbackslash newVec\{myNewVecB\}\{0.5, 1, 2\}
% Example: \textbackslash divVVec\{\textbackslash myNewVecA\}\{\textbackslash myNewVecB\} \(\xrightarrow{} \begin{pmatrix}1\\1\\1\end{pmatrix}\)
% #1 out, var name to save the result (vector) to
% #2 vector, the vector to divide
% #3 vectors, the list of vectors to divide with
\newcommand\divVVec[3][divVecResult]{
    \def\vecList{\xintCSVtoList{#3}}
    \getVecHeight{#2}
    \edef\sequence{\xintSeq[+1]{+1}{\vecHeightResult}}%
    % Iterate each row / dimension
    \xintFor* ##1 in \sequence \do {%
        % Assign value of start vector (param #2)
        \FPset\divValue{\xintNthElt{##1}{#2}}%
        % If this is the first iteration create a new vector
        \xintifForFirst{%
            % Iterate all vectors that are part of param #3
            \xintFor* ##2 in \vecList \do {%
                \FPifeq{\xintNthElt{##1}{##2}}{0}
                    \PackageError{glmatrix package}{In divVec function: division by zero}{you are trying to divide by 0}
                \else
                \fi
                % Divide with row value of current vector
                \FPdiv\divValue\divValue{\xintNthElt{##1}{##2}}%
            }
            \newVec{#1}{\divValue}
        }{%
            \xintFor* ##2 in \vecList \do {%
                \FPifeq{\xintNthElt{##1}{##2}}{0}
                    \PackageError{glmatrix package}{In divVec function: division by zero}{you are trying to divide by 0}
                \else
                \fi
                \FPdiv\divValue\divValue{\xintNthElt{##1}{##2}}%
            }
            \appendToVec{#1}{\csname #1\endcsname}{\divValue}
        }
    }
}

% Multiply vector by skalar.
% Equation: \(vector * scalar\)
% Example: \textbackslash newVec\{myNewVecA\}\{0.5, 1, 2\}
% Example: \textbackslash newScalar\{myNewScalar\}\{2\}
% Example: \textbackslash scaleVec\{\textbackslash myNewVecA\}\{\textbackslash myNewScalar\}
% Example: \textbackslash scaleVecResult \(\xrightarrow{} \begin{pmatrix}1\\2\\4\end{pmatrix}\)
% #1 out, var name to save the result (vector) to
% #2 vector, the vector to scale
% #3 scalar, the scalar to scale with
\newcommand\scaleVec[3][scaleVecResult] {%
    \xintFor* ##1 in #2 \do {%
        \FPmul\tmpVal{##1}{#3}
        \xintifForFirst{%
            \newVec{#1}{\tmpVal}
        }{%
            \appendToVec{#1}{\csname #1\endcsname}{\tmpVal}
        }
    }
}

% Divide vector by skalar.
% Equation: \(\frac{vector}{scalar}\)
% Example: \textbackslash newVec\{myNewVecA\}\{0.5, 1, 2\}
% Example: \textbackslash newScalar\{myNewScalar\}\{2\}
% Example: \textbackslash divVec\{\textbackslash myNewVecA\}\{\textbackslash myNewScalar\}
% Example: \textbackslash divVecResult \(\xrightarrow{} \begin{pmatrix}0.25\\0.5\\1\end{pmatrix}\)
% #1 out, var name to save the result (vector) to
% #2 vector, the vector to divide
% #3 scalar, the scalar to divide with
% TODO adjust function name, for function which divides vec with vec, negative scaling / shrinking
\newcommand\divVec[3][divVecResult] {%
    \FPifeq{#3}{0}
        \PackageError{glmatrix package}{In div vector by scalar}{divisor is 0}
    \else
    \fi
    \xintFor* ##1 in #2 \do {%
        \FPdiv\divNumRes{##1}{#3}
        \xintifForFirst{%
            \newVec{#1}{\divNumRes}
        }{%
            \appendToVec{#1}{\csname #1\endcsname}{\divNumRes}
        }
    }
}

% TODO negateVec vs negVec
% Negate the value of the given vector
% Example: \textbackslash newVec\{myNewVecA\}\{0.5, 1, 2\}
% Example: \textbackslash negateVecResult\{\textbackslash myNewVecA\}
% Example: \textbackslash negateVecResult \(\xrightarrow{} \begin{pmatrix}-0.5\\-1\\-2\end{pmatrix}\)
% #1 out, var name to save the result (vector) to
% #2 vector, the vector to negate
\newcommand\negateVec[2][negateVecResult]{
    \xintFor* ##1 in #2 \do {%
        \FPneg\tmp{##1}
        \xintifForFirst{%
            \newVec{#1}{\tmp}
        }{%
            \appendToVec{#1}{\csname #1\endcsname}{\tmp}
        }
    }
}

% Adds two vectors after multiplying the second one by a value.
% Equation: \(#1 = #2 + #3 * #4\)
% Example: \textbackslash newVec\{myNewVecA\}\{0.5, 1, 2\}
% Example: \textbackslash newVec\{myNewVecB\}\{0.5, 1, 2\}
% Example: \textbackslash newScalar\{myNewScalar\}\{2\}
% Example: \textbackslash scaleAndAddVec\{\textbackslash myNewVecA\}\{\textbackslash myNewVecB\}\{\textbackslash myNewScalar\}
% Example: \textbackslash scaleAndAddVecResult \(\xrightarrow{} \begin{pmatrix}1.5\\3\\6\end{pmatrix}\)
% #1 out, var name to save the result (vector) to
% #2 vector 1, the vector to add to
% #3 vector 2, the vector that is multiplied and then added
% #4 scalar, the scalar that is the muliplier
\newcommand\scaleAndAddVec[4][scaleAndAddVecResult]{
    \scaleVec{#3}{#4}
    \addVec[#1]{#2, \scaleVecResult}
}

% Lineraly interpolate between values of two vectors.
% Equation: \(out = (1-t) * vector1 + t * vector2\)
% Example: \textbackslash newVec\{myNewVecA\}\{0.5, 1, 2\}
% Example: \textbackslash newVec\{myNewVecB\}\{1, 2, 4\}
% Example: \textbackslash newScalar\{myNewScalar\}\{0.5\}
% Example: \textbackslash lerpVec\{\textbackslash myNewVecA\}\{\textbackslash myNewVecB\}\{\textbackslash myNewScalar\}
% Example: \textbackslash lerpVecResult \(\xrightarrow{} \begin{pmatrix}0.75\\1.5\\3\end{pmatrix}\)
% #1 out, var name to save the result (vector) to
% #2 vector 1, the start vector
% #3 vector 2, the end vector
% #4 t, interpolation amount range [0-1]
\newcommand\lerpVec[4][lerpVecResult]{
    % \def\vecList{\xintCSVtoList{#3}}
    \getVecHeight{#2}
    \edef\sequence{\xintSeq[+1]{+1}{\vecHeightResult}}%
    % Iterate each row / dimension
    \xintFor* ##1 in \sequence \do {%
        % Assign value of start vector (param #2)
        % \FPset\divValue{\xintNthElt{##1}{#2}}%
        % If this is the first iteration create a new vector
        \xintifForFirst{%
            \lerpScalar{\xintNthElt{##1}{#2}}{\xintNthElt{##1}{#3}}{#4}
            \newVec{#1}{\lerpScalarResult}
        }{%
            \lerpScalar{\xintNthElt{##1}{#2}}{\xintNthElt{##1}{#3}}{#4}
            \appendToVec{#1}{\csname #1\endcsname}{\lerpScalarResult}
        }
    }
}

% Performs a spherical linear interpolation between two vectors.
% Equation: TODO
% #1 out, var name to save the result (vector) to
% #2 vector 1, the start vector
% #3 vector 2, the end vector
% #4 t, interpolation amount range [0-1]
\newcommand\slerpVec[4][slerpVecResult] {%
    \angleVec{#2}{#3}
    \FPsin\angleTotal{\angleVecResult}
    \FPsub\ratioA{1}{#4}
    \FPmul\ratioA\ratioA{\angleVecResult}
    \FPsin\ratioA\ratioA
    \FPdiv\ratioA\ratioA\angleTotal
    \FPmul\ratioB{#4}{\angleVecResult}
    \FPsin\ratioB\ratioB
    \FPdiv\ratioB\ratioB\angleTotal
    \scaleVec[tmpA]{#2}{\ratioA}
    \scaleVec[tmpB]{#3}{\ratioB}
    \addVec[#1]{\tmpA, \tmpB}
}

% Perform a hermite interpolation with two control points.
% Equation: TODO
% #1 out, var name to save the result (vector) to
% #2 vector 1, the start vector
% #3 control 1, the first control point
% #4 control 2, the secont control point
% #5 vector 2, the end vector
% #6 t, interpolation amount range [0-1]
% TODO no idea if this works as intended
\newcommand\hermiteVec[6][hermiteVecResult] {%
    \FPmul\tQuad{#6}{#6}
    \FPmul\factorA{#6}{2}
    \FPsub\factorA\factorA{3}
    \FPmul\factorA\factorA\tQuad
    \FPadd\factorA\factorA{1}
    \FPsub\factorB{#6}{2}
    \FPmul\factorB\factorB\tQuad
    \FPadd\factorB\factorB{#6}
    \FPsub\factorC{#6}{1}
    \FPmul\factorC\factorC\tQuad
    \FPmul\factorD{#6}{2}
    \FPsub\factorD{3}\factorD
    \FPmul\factorD\factorD\tQuad
    \scaleVec[tmpA]{#2}{\factorA}
    \scaleVec[tmpB]{#3}{\factorB}
    \scaleVec[tmpC]{#4}{\factorC}
    \scaleVec[tmpD]{#5}{\factorD}
    \addVec[#1]{\tmpA, \tmpB, \tmpC, \tmpD}
}

% TODO
% \newcommand\bezier[3][distVecResult] {%
% }

% Calculate the skalarproduct/dotproduct of the given two vectors.
% Equation: \(X(v[1]) * X(v[2]) + Y(v[1]) * Y(v[2]) + ...\)
% Error: When both vectors do not have the same length
% Example: \textbackslash newVec\{myNewVecA\}\{0.5, 1, 2\}
% Example: \textbackslash newVec\{myNewVecB\}\{1, 2, 4\}
% Example: \textbackslash dotVec\{\textbackslash myNewVecA\}\{\textbackslash myNewVecB\}
% Example: \textbackslash dotVecResult \(\xrightarrow{} 10.5\)
% #1 out, var name to save the result (scalar) to
% #2 vector 1, the first vector
% #3 vector 2, the second vector
\newcommand\dotVec[3][dotVecResult] {%
    \FPifeq{\expandafter\xintLength\expandafter{#2}}{\expandafter\xintLength\expandafter{#3}}
    \else
        \PackageError{glmatrix package}{In dot product}{vector A and B do not have same length}
    \fi
    \FPset\res{0}
    \FPset\tmp{0}
    \edef\innersequence{\xintSeq[+1]{+1}{\expandafter\xintLength\expandafter{#2}}}%
    \xintFor* ##1 in \innersequence \do {%
        \FPset\valueA{\xintNthElt{##1}{#2}}
        \FPset\valueB{\xintNthElt{##1}{#3}}
        \FPmul\tmp\valueA\valueB
        \FPadd\res\res\tmp
    }
    \newScalar{#1}{\res}
}

% Calculate the length of the given vector.
% Equation: \(\sqrt{X(v)^2 + Y(v)^2 + ...}\)
% Example: \textbackslash newVec\{myNewVecA\}\{0.5, 1, 2\}
% Example: \textbackslash lenVec\{\textbackslash myNewVecA\}
% Example: \textbackslash lenVecResult \(\xrightarrow{} TODO scalar\)
% #1 out, var name to save the result (scalar) to
% #2 vector, the vector for which to calculate the length
\newcommand\lenVec[2][lenVecResult] {%
    \FPset\normTotal{0}
    \FPset\iterVal{0}
    % Iterate each value, multiply it with itself and sum up
    \xintFor* ##1 in #2 \do {%
        % No idea why this needs to be abs,
        % but negative values are a problem here
        \FPabs\iterVal{##1}
        \FPpow\tmpNormB\iterVal{2}
        \FPadd\normTotal\normTotal\tmpNormB
    }
    % Evaluate root for numbered result
    \FProot\normRooted\normTotal{2}
    % Set output var
    \newScalar{#1}{\normRooted}
}

% Calculate the squared length of the given vector.
% Equation: \(X(v)^2 + Y(v)^2 + ...\)
% Example: \textbackslash newVec\{myNewVecA\}\{0.5, 1, 2\}
% Example: \textbackslash sqrLenVec\{\textbackslash myNewVecA\}
% Example: \textbackslash sqrLenVecResult \(\xrightarrow{} TODO scalar\)
% #1 out, var name to save the result (scalar) to
% #2 vector, the vector for which to calculate the squared length
\newcommand\sqrLenVec[2][sqrLenVecResult] {%
    \FPset\normTotal{0}
    \FPset\iterVal{0}
    % Iterate each value, multiply it with itself and sum up
    \xintFor* ##1 in #2 \do {%
        % No idea why this needs to be abs,
        % but negative values are a problem here
        \FPabs\iterVal{##1}
        \FPpow\tmpNormB\iterVal{2}
        \FPadd\normTotal\normTotal\tmpNormB
    }
    % Set output var
    \newScalar{#1}{\normTotal}
}

% Divide all componendt of the vector by its length.
% Equation: \(out = \begin{pmatrix} \frac{X(v)}{\norm{v}}\\ \frac{Y(v)}{\norm{v}}\\ ...\end{pmatrix}\)
% Example: \textbackslash newVec\{myNewVecA\}\{0.5, 1, 2\}
% Example: \textbackslash normalizeVec\{\textbackslash myNewVecA\}
% Example: \textbackslash normalizeVecResult \(\xrightarrow{} TODO vector\)
% #1 out, var name to save the result (vector) to
% #2 vector, the vector which to normalize
\newcommand\normalizeVec[2][normalizeVecResult] {%
    \lenVec{#2}
    \divVec[#1]{#2}{\lenVecResult}
}

% TODO this only works for 3d vectors currently
% Calculate crossproduct for the two given vectors.
% Equation: \(out = \begin{pmatrix} ab*bc - ac*bb\\ ac*ba-aa*bc\\ aa*bb-ab*ba \end{pmatrix}\)
% Error: When one of both vectors does not have length 3
% Example: \textbackslash newVec\{myNewVecA\}\{0.5, 1, 2\}
% Example: \textbackslash newVec\{myNewVecB\}\{0.5, 1, 2\}
% Example: \textbackslash crossVec\{\textbackslash myNewVecA\}\{\textbackslash myNewVecB\}
% Example: \textbackslash crossVecResult \(\xrightarrow{} TODO vector\)
% #1 out, var name to save the result (vector) to
% #2 vector, the first vector
% #3 vector, the second vector
\newcommand\crossVec[3][crossVecResult] {%
    \FPifeq{\expandafter\xintLength\expandafter{#2}}{3}
    \else
        \PackageError{glmatrix package}{In cross product}{vector A is not of length 3}
    \fi
    \FPifeq{\expandafter\xintLength\expandafter{#3}}{3}
    \else
        \PackageError{glmatrix package}{In cross product}{vector A is not of length 3}
    \fi
    \FPset\aa{\xintNthElt{1}{#2}}
    \FPset\ab{\xintNthElt{2}{#2}}
    \FPset\ac{\xintNthElt{3}{#2}}
    \FPset\ba{\xintNthElt{1}{#3}}
    \FPset\bb{\xintNthElt{2}{#3}}
    \FPset\bc{\xintNthElt{3}{#3}}
    \FPset\tmp{0}
    % x component
    \FPset\scalarX{0}
    \FPmul\scalarX\ab\bc
    \FPmul\tmp\ac\bb
    \FPsub\scalarX\scalarX\tmp
    % y component
    \FPset\scalarY{0}
    \FPmul\scalarY\ac\ba
    \FPmul\tmp\aa\bc
    \FPsub\scalarY\scalarY\tmp
    % z component
    \FPset\scalarZ{0}
    \FPmul\scalarZ\aa\bb
    \FPmul\tmp\ab\ba
    \FPsub\scalarZ\scalarZ\tmp
    \newVec{#1}{\scalarX,\scalarY,\scalarZ}
}

% Get the angle between the given two vectors in radian.
% Equation: \(out = angle(v1, v2)\)
% Example: \textbackslash newVec\{myNewVecA\}\{0.5, 1, 2\}
% Example: \textbackslash newVec\{myNewVecB\}\{0.5, 1, 2\}
% Example: \textbackslash angleVec\{\textbackslash myNewVecA\}\{\textbackslash myNewVecB\}
% Example: \textbackslash angleVecResult \(\xrightarrow{} TODO scalar\)
% #1 out, var name to save the result (scalar) to
% #2 vector 1, the first vector
% #3 vector 2, the second vector
% TODO 180 degree does not work
% TODO check if vector length to short (1)
\newcommand\angleVec[3][angleVecResult] {%
    \dotVec[dotVW]{#2}{#3}
    \lenVec[normV]{#2}
    \lenVec[normW]{#3}
    \mulScalar[mulNormVW]{\normV}{\normW}
    \divScalar{\dotVW}{\mulNormVW}
    \arccosScalar[#1]{\scalarDivResult}
}

% TODO this does only work in 2d
% \newcommand\signedAngleVec[3][distVecResult] {%
% }

% Calculates the euclidian distance between the given two vectors.
% Equation: TODO
% Error: When vector 1 and 2 do not have the same length
% Example: \textbackslash newVec\{myNewVecA\}\{0.5, 1, 2\}
% Example: \textbackslash newVec\{myNewVecB\}\{0.5, 1, 2\}
% Example: \textbackslash distVec\{\textbackslash myNewVecA\}\{\textbackslash myNewVecB\}
% Example: \textbackslash distVecResult \(\xrightarrow{} TODO scalar\)
% #1 out, var name to save the result (scalar) to
% #2 vector 1, the first vector
% #3 vector 2, the second vector
\newcommand\distVec[3][distVecResult] {%
    \sqrDistVec{#2}{#3}
    \FProot\res\sqrDistVecResult{2}
    \newScalar{#1}{\res}
}

% Calculates the squared euclidian distance between two vectors.
% Error: When vector 1 and 2 do not have the same length
% Example: \textbackslash newVec\{myNewVecA\}\{0.5, 1, 2\}
% Example: \textbackslash newVec\{myNewVecB\}\{0.5, 1, 2\}
% Example: \textbackslash distVec\{\textbackslash myNewVecA\}\{\textbackslash myNewVecB\}
% Example: \textbackslash distVecResult \(\xrightarrow{} TODO scalar\)
% #1 out, var name to save the result (scalar) to
% #2 vector 1, the first vector
% #3 vector 2, the second vector
\newcommand\sqrDistVec[3][sqrDistVecResult] {%
    \getVecHeight[getVecHeightA]{#2}
    \getVecHeight[getVecHeightB]{#3}
    \FPifeq{\getVecHeightA}{\getVecHeightB}
    \else
        \PackageError{glmatrix package}{In sqrt distance}{vector #2 and #3 do not have same length}
    \fi
    \FPset\res{0}
    \FPset\tmp{0}
    \edef\sequence{\xintSeq[+1]{+1}{\getVecHeightA}}%
    \xintFor* ##1 in \sequence \do {%
        \FPset\valueA{\xintNthElt{##1}{#2}}
        \FPset\valueB{\xintNthElt{##1}{#3}}
        \FPsub\tmp\valueA\valueB
        \FPmul\tmp\tmp\tmp
        \FPadd\res\res\tmp
    }
    \newScalar{#1}{\res}
}

% Multiply the given vector with a matrix (for transformation).
% Error: When matrix and vector have a dimension mismatch
% Example: \textbackslash newVec\{myNewVec\}\{0.5, 1, 2\}
% Example: \textbackslash newMatrix\{myNewMatrix\}\{\{0.5, 1, 2\}\{0.5, 1, 2\}\{0.5, 1, 2\}\}
% Example: \textbackslash transformVec\{\textbackslash myNewMatrix\}\{\textbackslash myNewVec\}
% Example: \textbackslash transformVecResult \(\xrightarrow{} TODO vector\)
% #1 out, var name to save the result (vector) to
% #2 matrix, the matrix to transform the vetor with
% #3 vector, the vector to transform
% TODO if the vector is smaller than the matrix the missing components, should automatically be set to one
\newcommand\transformVec[3][transformVecResult] {%
    \getVecHeight{#3}%
    \getMatHeight{#2}%
    \getMatWidth{#2}%
    % Check that matrix and vector can be multiplied
    \FPifeq{\matWidthResult}{\vecHeightResult}
    \else
        \PackageError{glmatrix package}{Dimension mismatch for matrix and vector, could not multiply}{width (\matWidthResult) of matrix #2 and height (\vecHeightResult) of vector #3 do not match}
    \fi
    % Create iterator sequences from height of vector and matrix
    \edef\vecHeightSequence{\xintSeq[+1]{+1}{\vecHeightResult}}%
    \edef\matHeightSequence{\xintSeq[+1]{+1}{\matHeightResult}}%
    % Iterate all rows of the matrix
    \xintFor* ##2 in \matHeightSequence \do {%
        \FPset\cellVal{0}
        % Differentiate between first and other entries to create the new vector variable
        \xintifForFirst{%
            % Now iterate each entry of the current row
            % Use this indec to access the entry in the matrix and in the vector
            \xintFor* ##3 in \vecHeightSequence \do {%
                \FPmul\tmp{%
                    \xintNthElt{##3}{\xintNthElt{##2}{#2}}%
                }{%
                    \xintNthElt{##3}{#3}%
                }
                % Add to storage variable
                \FPadd\cellVal\cellVal\tmp
            }
            \FPclip\cellVal\cellVal
            \newVec{#1}{\cellVal}
        }{
            \xintFor* ##3 in \vecHeightSequence \do {%
                \FPmul\tmp{%
                    \xintNthElt{##3}{\xintNthElt{##2}{#2}}%
                }{%
                    \xintNthElt{##3}{#3}%
                }
                \FPadd\cellVal\cellVal\tmp
            }
            \FPclip\cellVal\cellVal
            \appendToVec{#1}{\csname #1\endcsname}{\cellVal}
        }
    }
    % Test if width and height is correct
    \expandafter\getVecHeight\expandafter{\csname #1\endcsname}%
    \FPifeq{\matHeightResult}{\vecHeightResult}
    \else
        \PackageError{glmatrix package}{Vector height does not match expectation}{height of vector is \vecHeightResult but should be \matHeightResult}
    \fi
}

% Rotate the vector by #4 degrees around the X axis.
% #1 out, var name to save the result (vector) to
% #2 a, vector value
% #3 b, the origin of the rotation
% #4 c, the rotation angle in radians
% TODO needs atleast 3 components
% \newcommand\rotateX[4][rotateXVecResult] {%
% }

% Rotate the vector by #4 degrees around the Y axis.
% #1 out, var name to save the result (vector) to
% #2 a, vector value
% #3 b, the origin of the rotation
% #4 c, the rotation angle in radians
% TODO needs atleast 3 components
% \newcommand\rotateY[3][rotateYVecResult] {%
% }

% Rotate the vector by #4 degrees around the Z axis.
% #1 out, var name to save the result (vector) to
% #2 a, vector value
% #3 b, the origin of the rotation
% #4 c, the rotation angle in radians
% TODO does also work for 2d vectors
% \newcommand\rotateZ[3][rotateZVecResult] {%
% }

% TODO
% \newcommand\transformQuat[3][distVecResult] {%
% }

% TODO BIG
% Multiply a vector (column) with itself (row) resulting in a matrix.
% Error: When matrix and vector have a dimension mismatch
% #1 out, var name to save the result (matrix) to
% #2 vector, the vector to multiply with itself
\newcommand\mulVecAsColAndRow[2][mulMatsResult]{%
    \newMatRow{vecA}{\xintNthElt{1}{#2}}
    \newMatRow{vecB}{\xintNthElt{2}{#2}}
    \newMatRow{vecC}{\xintNthElt{3}{#2}}
    \newMatFromRowVec{tmpMatA}{\vecA, \vecB, \vecC}
    \vecToMatRow{row}{#2}
    \newMat{tmpMatB}{\row}
    \mulMat[#1]{\tmpMatA}{\tmpMatB}
}

%%%%%%%%%%%%%%%%%%%%%
% Vector Comparison %
%%%%%%%%%%%%%%%%%%%%%
% Returns whether or not the vectors have the same values. 
% \newcommand\equalsVec[3][equalsVecResult]{
% }

% Return the componentwise minimum of the given two vectors.
% Equation: \(\begin{pmatrix} \min(X(v1), X(v2))\\ \min(Y(v1), Y(v2))\\ ... \end{pmatrix}\)
% Error: When vector 1 and 2 do not have the same length
% #1 out, var name to save the result (vector) to
% #2 vector 1, the first vector
% #3 vector 2, the second vector
\newcommand\minVec[3][minVecResult]{
    \FPifeq{\expandafter\xintLength\expandafter{#2}}{\expandafter\xintLength\expandafter{#3}}
    \else
        \PackageError{glmatrix package}{In vector minimum}{vector A and B do not have same length}
    \fi
    \edef\innersequence{\xintSeq[+1]{+1}{\expandafter\xintLength\expandafter{#2}}}%
    \xintFor* ##1 in \innersequence \do {%
        \minScalar{\xintNthElt{##1}{#2}}{\xintNthElt{##1}{#3}}
        \xintifForFirst{%
            \newVec{#1}{\minScalarResult}
        }{%
            \appendToVec{#1}{\csname #1\endcsname}{\minScalarResult}
        }
    }
}

% Return the componentwise maximum of the given two vectors.
% Equation: \(\begin{pmatrix} \max(X(v1), X(v2))\\ \max(Y(v1), Y(v2))\\ ... \end{pmatrix}\)
% Error: When vector 1 and 2 do not have the same length
% #1 out, var name to save the result (vector) to
% #2 vector 1, the first vector
% #3 vector 2, the second vector
\newcommand\maxVec[3][maxVecResult]{
    \FPifeq{\expandafter\xintLength\expandafter{#2}}{\expandafter\xintLength\expandafter{#3}}
    \else
        \PackageError{glmatrix package}{In vector minimum}{vector A and B do not have same length}
    \fi
    \edef\innersequence{\xintSeq[+1]{+1}{\expandafter\xintLength\expandafter{#2}}}%
    \xintFor* ##1 in \innersequence \do {%
        \maxScalar{\xintNthElt{##1}{#2}}{\xintNthElt{##1}{#3}}
        \xintifForFirst{%
            \newVec{#1}{\maxScalarResult}
        }{%
            \appendToVec{#1}{\csname #1\endcsname}{\maxScalarResult}
        }
    }
}

%%%%%%%%%%%%%
% Functions %
%%%%%%%%%%%%%
% Apply the given function to a list of variables.
% TODO try to use a string as #2 instead of the integer, 
% this would allow for dynamic function calls, not requireing the if else
% utelise \csname #2\endcsname for that
% #1 list of commands (var names) that should be altered
% #2 the chosen function
% #3 params for function
% \newcommand\forEach{
% }

% TODO
% #1 a, a
% #2 a, a
% #3 a, a
\newcommand\solveVecEQL[3][solveVecEQLResult]{
    \getVecHeight{#3}%
    \edef\vecHeightSequence{\xintSeq[+1]{+1}{\vecHeightResult}}%
    \xintFor* ##1 in \vecHeightSequence \do {%
        \FPneg\tmpVar{\xintNthElt{##1}{#3}}
        \FPlsolve\glCurSol{\xintNthElt{##1}{#2}}{\tmpVar}
        % \maxScalar{\xintNthElt{##1}{#2}}{\xintNthElt{##1}{#3}}
        % \xintifForFirst{%
        %     \newVec{#1}{\maxScalarResult}
        % }{%
        %     \appendToVec{#1}{\csname #1\endcsname}{\maxScalarResult}
        % }
        \xintifForFirst{%   
            \setScalar{#1}{\glCurSol}
        }{%
            \FPifeq{\csname #1\endcsname}{\glCurSol}
            \else
                \PackageError{glmatrix package}{Could not solve equation}{varying results for different dimensions}
            \fi
        }
    }
}